\documentclass{situationist-preprint}

\title{Meaning Without God}
\msid{EXI-11-a}
\msfamily{EXI-11}
\msdiscipline{Meaning, Death and the Absurd}
\mstype{Principal manuscript}
\mskeywords{meaning in life; God; cosmic purpose; purpose theory; supernaturalism; naturalism; Wolf; Metz; Cottingham; Nozick}

\begin{document}
\maketitle

\begin{abstract}
Can life be meaningful without God? Supernaturalists hold that meaning requires a purpose assigned to human life by its creator, or a relation to an infinite and imperishable good, and that without these every human pursuit comes to nothing. We argue that life can be meaningful without God, because meaning does not require a cosmic purpose, and that a theistic framework would nevertheless add to it. The first claim rests on the fitting-fulfilment account of meaning: a life is meaningful to the extent that its subject actively engages with what is worthwhile, and the worth of knowledge, of persons and of beauty is not the worth of a role assigned to us by a designer. Purpose assigned by another is neither necessary for meaning nor sufficient for it. The second claim rests on the observation that meaning comes in degrees and grows with the worth of what one engages: a framework containing a being of maximal worth, a purpose that unifies one's projects, and the preservation of what is good offers objects and relations of engagement that a naturalistic world lacks. We read the question as one about dependence, so that the answer is independent of whether God exists.
\end{abstract}

\section{The problem}

Tolstoy recounts how, at the height of his success, he was stopped by the question of what would come of everything he did, and found that if everything ended in death there was no answer \citep{Tolstoy1882}. Only faith, which relates the finite to the infinite, gave him one. The argument he found in his own crisis has been stated by philosophers in several forms. On purpose theories, a life is meaningful insofar as it fulfils a purpose that God has assigned to it; without God there is no such purpose, and so no meaning \citep{Craig2008}. On theories of the infinite, meaning requires a connection to something of unlimited worth, and only God is such a thing \citep[ch.~6]{Nozick1981}. On \citeauthor{Cottingham2003}'s more cautious view, the pursuit of meaning requires a moral framework and a ground for hope that the good will prevail, and the religious outlook provides these in a way that its rivals do not \citep{Cottingham2003}.

Naturalists reply that meaning is available in a world without God. On subjectivist views, meaning is a matter of what a person cares about \citep{Taylor1970}. On objective naturalist views, it is a matter of standing in the right relation to values that are features of the natural world or of what rational agents would endorse. \citet{Wolf2010} defends a hybrid: meaning arises when the subject's active engagement meets the objective worth of its object. \citet{Metz2013} surveys the supernaturalist theories at length, rejects the purpose theory and the soul-centered theories as accounts of what meaning requires, and defends a naturalist theory on which meaning is a matter of orienting one's rational nature toward fundamental conditions of human existence. Metz also distinguishes extreme supernaturalism, on which God is necessary for any meaning, from moderate supernaturalism, on which God is necessary only for a greater or ultimate meaning.

We defend that position. Section~2 sets out the framework, section~3 argues that meaning does not require cosmic purpose, and section~4 argues that a theistic framework would add to meaning. Section~5 answers objections.

\paragraph{The question.} REL-01 argues that a maximally great being exists necessarily. If that is right, there is no possible world without God, and the question whether life would be meaningful in such a world is a counterpossible whose answer is unclear. We therefore read the question, as most of its participants do, as a question about dependence: whether the conditions that make a life meaningful include God's existence, a relation to God, or a purpose assigned by God. So read, the question has a determinate answer whether or not REL-01 is correct, and the answer also settles a question of evident practical interest, namely whether the life of someone who does not believe in God and does not relate to God can be meaningful.

\section{Framework}

\begin{definition}[Meaning]\label{d:meaning}
A life is \emph{meaningful} to the extent that its subject is actively engaged with what is worthwhile.
\end{definition}

Definition~\ref{d:meaning} is the fitting-fulfilment account, defended in EXI-17, which also argues that achievement is not required. EXI-01 adds that meaning so understood is available to everyone.

\begin{definition}[Cosmic purpose]\label{d:purpose}
A life has a \emph{cosmic purpose} if a creator of the universe intends that it play a certain role in a plan for the universe as a whole. A life \emph{fulfils} its cosmic purpose if it plays that role.
\end{definition}

\begin{definition}[Purpose theory]\label{d:pt}
The \emph{purpose theory} holds that a life is meaningful only if it has a cosmic purpose and fulfils it.
\end{definition}

\begin{principle}[Worth is not conferred by assignment]\label{p:worth}
That an activity has been assigned to an agent as a purpose by another, however powerful, does not by itself make the activity worthwhile.
\end{principle}

Principle~\ref{p:worth} follows from the account of the good in ETH-01, on which the good is what idealized reflective endorsement converges on, and it is consistent with the account of moral grounding in MTE-20, on which divine commands confirm moral obligations rather than create them. The principle does not deny that God may value what is valuable, or command what is good; it denies that the assignment is what makes it good.

\begin{principle}[Degree]\label{p:degree}
Other things equal, engagement with what is of greater worth, or engagement that relates one to more of what is worthwhile, makes a life more meaningful.
\end{principle}

Principle~\ref{p:degree} is implicit in Definition~\ref{d:meaning}, which makes meaning a matter of extent. It also expresses the insight in \citeauthor{Nozick1981}'s linkage view: a life gains meaning by connection with value beyond itself, and more by connection with greater value \citep[ch.~6]{Nozick1981}.

\section{Meaning without cosmic purpose}

\begin{lemma}[Purpose is not sufficient]\label{l:notsuff}
Fulfilling a cosmic purpose is not sufficient for a meaningful life.
\end{lemma}

\begin{proof}
Consider a creator who brings human beings into existence so that they may serve as food for another species, or as spectacles for its amusement, and suppose a life that fulfils that purpose. It has a cosmic purpose and fulfils it. By Principle~\ref{p:worth}, the assignment does not make serving as food worthwhile, and the subject need not be engaged with anything worthwhile in living it. By Definition~\ref{d:meaning}, such a life is not made meaningful by its purpose. Reflection agrees: a purpose of this kind degrades the life it is given to.
\end{proof}

The supernaturalist will reply that God's purposes are good. That is right, but it concedes the point of Lemma~\ref{l:notsuff}: what makes the fulfilment of a cosmic purpose meaningful, on that reply, is the goodness of the purpose, not the fact that it is assigned.

\begin{lemma}[Worth without purpose]\label{l:worth}
There are activities that are worthwhile independently of whether anyone has assigned them as a cosmic purpose.
\end{lemma}

\begin{proof}
Coming to understand the structure of the natural world, caring for a sick child, and making or attending to a beautiful thing realize knowledge, well-being, relationships and beauty, which ETH-01 identifies as components of the good. By Principle~\ref{p:worth}, the worth of these activities is not conferred by their assignment; and ETH-01's account of the good, as what idealized reflective endorsement converges on, does not make it depend on any assignment. So they are worthwhile whether or not they are assigned as cosmic purposes.
\end{proof}

\begin{theorem}\label{t:main}
Meaning does not require a cosmic purpose. A life that has no cosmic purpose, or does not fulfil the one it has, can be meaningful.
\end{theorem}

\begin{proof}
By Lemma~\ref{l:worth}, some activities are worthwhile independently of any assignment. A subject who is actively engaged with such activities satisfies Definition~\ref{d:meaning}, whether or not her life has a cosmic purpose and whether or not she fulfils it. Her life is therefore meaningful to that extent. The purpose theory of Definition~\ref{d:pt} is false.
\end{proof}

\begin{corollary}[The naturalist's life]\label{c:atheist}
The life of someone who does not believe in God and does not relate to God can be meaningful.
\end{corollary}

\begin{proof}
Definition~\ref{d:meaning} makes no reference to belief in God or relation to God. By Theorem~\ref{t:main}, such a person's engagement with what is worthwhile makes her life meaningful.
\end{proof}

\section{What theism would add}

\begin{proposition}[Theism adds]\label{p:adds}
A theistic framework, if true, would make available to a life more meaning than it could have in a world without God.
\end{proposition}

\begin{proof}
A theistic framework of the kind REL-01 defends contains a being of maximal worth. Engagement with that being, through love, worship, obedience and the attempt to understand it, would be engagement with the greatest object of worth there is; by Principle~\ref{p:degree}, it would make a life more meaningful than engagement with lesser objects, other things equal. The framework also offers relations to what is worthwhile that a naturalistic world lacks. It unifies a person's projects under a single purpose that is itself good, so that her separate engagements are connected with one another and with the good of the whole. And it holds out the preservation of what is good, so that the goods one serves are not finally lost. Each of these relates the life to more of what is worthwhile, and by Principle~\ref{p:degree} each would add to its meaning.
\end{proof}

\begin{remark}\label{r:status}
Proposition~\ref{p:adds} is conditional on the framework. If REL-01 is right, the additional meaning is actually available, and is available to those who engage with God; if REL-01 is wrong, it is not available, and Theorem~\ref{t:main} ensures that life is meaningful nonetheless. Our argument does not use REL-01 as a premise. The claim that the goods one serves are preserved should be held with the caution REL-08 recommends: survival of death is coherent, but no evidence available before death could confirm it.
\end{remark}

Theorem~\ref{t:main} and Proposition~\ref{p:adds} together give a position that is neither supernaturalist nor dismissive of the supernaturalist's insight. Meaning has a floor that does not depend on God, and theism, if true, would raise its ceiling. It agrees with extreme supernaturalism's critics that God is not necessary for meaning, and with moderate supernaturalism that God would make more meaning available \citep{Metz2013}; it does not follow the moderate view in making God necessary for any particular grade of meaning.

\section{Objections and replies}

\begin{objection}[Everything comes to nothing]\label{o:tolstoy}
Without God, everything one does will eventually be lost: one dies, one's works decay, and the universe ends. What comes to nothing was nothing \citep{Tolstoy1882,Craig2008}.
\end{objection}

\begin{reply}
The inference from transience to worthlessness is invalid. A friendship that ends with death was a friendship; the relief of a patient's pain was relief while it lasted. Worth realized at a time is not cancelled by later loss, and the claim that it is would make the worth of a good depend on what happens after it, which neither ETH-01's account of the good nor any plausible account allows. EXI-19 argues that death, far from voiding a life, gives it a shape by limiting it. What permanence would add is accommodated by Proposition~\ref{p:adds}: preservation of the good would add meaning; its absence does not remove the meaning already there.
\end{reply}

\begin{objection}[The regress of linkage]\label{o:nozick}
On the linkage view, a finite life has meaning by connection with something larger. That larger thing has meaning only if it is connected with something larger still, and the regress must end in something unlimited, which only God can be \citep[ch.~6]{Nozick1981}.
\end{objection}

\begin{reply}
The regress arises only if worth is always borrowed from something else. On our account, the regress ends in things that are worthwhile in themselves, such as knowledge, well-being and beauty (ETH-01). A life linked to these is meaningful without further linkage. The demand for an unlimited terminus expresses a wish for meaning without limit, not a condition of meaning as such. Proposition~\ref{p:adds} grants that an unlimited object would add meaning; it does not follow that without it there is none.
\end{reply}

\begin{objection}[The framework of hope]\label{o:cottingham}
The pursuit of meaning is not a solitary engagement with isolated goods. It requires a stable moral framework and a ground for hope that our efforts toward the good are not in vain. Without God, these are fragile, and so is meaning \citep{Cottingham2003}.
\end{objection}

\begin{reply}
We accept much of this. A framework that unifies one's engagements and grounds hope that the good prevails is among the additions Proposition~\ref{p:adds} identifies. But the fragility the objection describes affects the security of meaning and the reasonableness of certain hopes (EXI-15), not its existence. The engaged and worthwhile life of a person without such a framework is not made meaningless by its exposure to discouragement.
\end{reply}

\begin{objection}[Purpose is what meaning means]\label{o:concept}
The word \enquote{meaning} in \enquote{the meaning of life} just asks for a purpose. A life without a purpose assigned to it has no meaning by definition, and the naturalist has changed the subject.
\end{objection}

\begin{reply}
Usage does not support the claim. The question of what makes a life meaningful is asked and answered by believers and unbelievers alike in terms of what a life is devoted to and whether that is worthy of devotion. Lemma~\ref{l:notsuff} shows that purpose in the sense of assignment does not capture this, since a purpose may be assigned that makes a life worse. If the purpose theorist replies that only good purposes count, the notion she uses is no longer purpose as such but engagement with the good, which is Definition~\ref{d:meaning}.
\end{reply}

\begin{objection}[Instability]\label{o:instab}
If a life can be fully meaningful without God, nothing is left for theism to add; if theism adds, the life without God is deficient in meaning.
\end{objection}

\begin{reply}
The dilemma assumes that meaning is all or nothing. On Definition~\ref{d:meaning} it comes in degrees. A life can be meaningful, and deeply so, while a further relation to a greater good would make it more meaningful still. That a life could be more meaningful is not a deficiency in any sense that threatens Theorem~\ref{t:main}; it is true of every finite life.
\end{reply}

\section{What the result does not establish}

The argument does not show that God exists or that God does not; REL-01 addresses that question, and our conclusion is independent of its answer. It does not show how much meaning theism would add, nor that the addition outweighs other considerations bearing on religious commitment. Nor does it show that every life without God is meaningful: meaning requires engagement with what is worthwhile, and lives with and without God can fail that condition alike. Finally, it relies on the objectivity of worth established elsewhere (ETH-01, MTE-20); a reader who rejected that would need a different route to Theorem~\ref{t:main}.

\section{Conclusion}

Life can be meaningful without God. Meaning consists in active engagement with what is worthwhile, and the worth of the objects of engagement does not derive from their being assigned to us as a cosmic purpose; assigned purpose is neither sufficient for meaning nor necessary for it. A theistic framework would add to meaning, by offering an object of maximal worth, a purpose that unifies one's engagements, and the preservation of the good. The floor of meaning is independent of God; the ceiling, if REL-01 is right, is raised by God.

\begin{thebibliography}{9}\small
\bibitem[Cottingham(2003)]{Cottingham2003} Cottingham, J. (2003). \emph{On the Meaning of Life}. Routledge.
\bibitem[Craig(2008)]{Craig2008} Craig, W. L. (2008). The absurdity of life without God. In \emph{Reasonable Faith: Christian Truth and Apologetics}, 3rd ed. Crossway.
\bibitem[Metz(2013)]{Metz2013} Metz, T. (2013). \emph{Meaning in Life: An Analytic Study}. Oxford University Press.
\bibitem[Nozick(1981)]{Nozick1981} Nozick, R. (1981). \emph{Philosophical Explanations}. Harvard University Press.
\bibitem[Taylor(1970)]{Taylor1970} Taylor, R. (1970). \emph{Good and Evil: A New Direction}. Macmillan.
\bibitem[Tolstoy(1882)]{Tolstoy1882} Tolstoy, L. (1882). \emph{A Confession}. Cited by chapter in any edition.
\bibitem[Wolf(2010)]{Wolf2010} Wolf, S. (2010). \emph{Meaning in Life and Why It Matters}. Princeton University Press.
\end{thebibliography}

\end{document}
